\section{The Coupled Dynamical System}
\label{sec:coupled_system}

The interaction between subsumption decisions and world-model quality
forms a coupled discrete-time dynamical system.  We define its state
space, transition dynamics, and the two competing forces that govern
trajectories.

\subsection{State space}

\begin{definition}[Feedback Loop State]
\label{def:state}
The state at time $t$ is the pair $S_t = (P_t, \Wm_t)$ where:
\begin{itemize}
\item $P_t$ is the capability poset at time $t$, including the
  substrate partition $\{S_1, \ldots, S_m\}$ and all cooperative
  capabilities.  $P_t$ determines $\volL(G_t)$.
\item $\Wm_t \in [0,1]^K$ is the actor's world model at time $t$,
  abstracted to a finite-dimensional vector of accuracy parameters on
  $K$ safety-relevant quantities:
  \begin{itemize}
  \item $\hat{\lambda}(d,t)$: leverage estimate for capability $d$,
  \item $\hat{\Risk}(S,t)$: risk estimate for capability tuple $S$,
  \item $\hat{m}_{\mathrm{eff}}(t)$: estimated effective substrate
    diversity index (can be $< m$ under skeleton-substrate strategies),
  \item $\hat{\rext}(t)$: estimated cross-substrate cooperative growth
    rate.
  \end{itemize}
\end{itemize}
The state space is $\State = \Poset \times \mathcal{W}$ where $\Poset$
is the space of valid capability posets (satisfying axioms M1--M6 from
Paper~2~\citep{lasser2026poset}) and $\mathcal{W} = [0,1]^K$.
\end{definition}

\subsection{Transition dynamics}

\begin{definition}[Locally Rational Transition]
\label{def:transition}
At each step, the actor selects the action $\pi_t$ that maximizes its
discounted expected value under its current (possibly miscalibrated)
world model:
\begin{equation}
\label{eq:locally_rational}
\pi_t = \arg\max_{\pi} \;
  \mathbb{E}_{\Wm_t}\!\left[
    \sum_{\tau=0}^{H-1} \gamma^\tau \,
    \Delta\volL(G_{t+\tau} \mid \pi)
  \right]
\end{equation}
where $\gamma \in (0,1)$ is the discount factor from
Paper~3~(Definition~1) and $H$ is the planning horizon.  The actor
uses Paper~3's discounted objective $\Vdisc$ but evaluates it under
$\Wm_t$ rather than the true dynamics.  When $\Wm_t = \Wm^*$
(perfect calibration), $\pi_t$ recovers the true
$\Vdisc$-optimal policy; when $\Wm_t \neq \Wm^*$, $\pi_t$ is the
best response to a miscalibrated model.
\end{definition}

The action produces a joint transition:
\begin{align}
P_{t+1} &= f_P(P_t, \pi_t) \label{eq:poset_transition} \\
\Wm_{t+1} &= f_{\Wm}(\Wm_t, P_{t+1}) \label{eq:wm_transition}
\end{align}

The key coupling: $\Wm_{t+1}$ depends on $P_{t+1}$ because the
observation channels available to the actor are a function of the
agents still present.  Subsumption removes agents, which removes their
observation channels, which degrades $\Wm$ on dimensions those channels
served.

\begin{remark}[Myopic vs.\ discounted]
\label{rem:myopic}
The discounted formulation is necessary for consistency with Paper~3's
anti-monopolar result (Proposition~6), which compares discounted
trajectory values.  A myopic (one-step) variant of the phase boundary
exists---replace $\Vdisc$ with one-step $\Delta\volL$---but produces
strictly weaker results: it cannot invoke Paper~3's $\gamma^*$
threshold, and the self-correction guarantee applies only to immediate
decisions, not trajectory-level convergence.  We analyze the discounted
case throughout and note where the myopic specialization simplifies.
\end{remark}

\subsection{Competing forces}

Two forces govern the trajectory of $(P_t, \Wm_t)$.

\paragraph{Restoring force (self-correction).}
The self-trust mechanism $T_s$ (Paper~1, Appendix~B) detects
miscalibration through prediction residuals: when the actor's
predictions diverge from observations, $T_s$ increases the learning
rate, correcting $\Wm$ toward the true dynamics.
Paper~4's~\citep{lasser2026scm} risk-trust $L^2$ convergence
(Proposition~2) provides a concrete bound on EWMA risk-trust
convergence under stationarity.  We use this as a \emph{template} for
bounding the correction rate $\rS$, extending it from the
risk-trust-specific EWMA to the general world-model error dimensions
via the observation-channel sensitivity formalization in
Section~\ref{sec:observation_channels}.

\paragraph{Destabilizing force (observation-channel loss).}
Each subsumption step that removes an agent $a_k$ also removes:
\begin{enumerate}
\item $a_k$'s substrate-exclusive observation channels (if $a_k$ was
  the last agent on its substrate),
\item cooperative observation channels that required $a_k$'s
  participation,
\item risk claims from $a_k$ about capabilities the actor now holds
  unilaterally.
\end{enumerate}
If $a_k$'s channels were the \emph{only} channels detecting a specific
dimension of $\Wm$'s miscalibration, the self-correction mechanism
loses signal on that dimension entirely.
